\booksection{Introduction}{Le matin ordinaire}

Un matin ordinaire, une personne seule ouvre son ordinateur. En quelques minutes, elle peut demander à une intelligence artificielle de traduire un document, d'analyser un marché, d'écrire un programme, de créer une présentation, d'expliquer un concept scientifique difficile, de bâtir un plan d'affaires ou encore de formuler des arguments contradictoires sur un même sujet. Il y a encore peu de temps, chacune de ces tâches aurait exigé des heures de travail, mobilisé plusieurs compétences différentes, ou fait appel à des spécialistes. Aujourd'hui, pour un nombre croissant d'individus, ce qui relevait hier de l'exceptionnel tend à devenir banal.

Ce glissement est d'une importance considérable. Nous avons l'habitude d'associer les grandes révolutions technologiques à des machines qui ont décuplé notre force physique : la machine à vapeur, l'électricité, l'automobile ou encore l'énergie nucléaire. Chacune d'elles a modifié notre rapport à l'espace, au temps ou à la matière. Ce qui se dessine aujourd'hui est d'une autre nature. Pour la première fois, ce n'est pas notre force qui est amplifiée, mais une partie de nos capacités cognitives.

Ce basculement soulève une question simple, mais d'une profondeur peu commune :

\begin{aiquote}
À quoi pourrait ressembler la société lorsque l'intelligence artificielle ne sera plus simplement un outil que nous utilisons, mais une infrastructure permanente de notre vie économique, professionnelle, éducative, culturelle et sociale ?
\end{aiquote}

La tentation est alors de se demander : «Que peut faire l'IA ?». Cette question, légitime, reste cependant trop étroite. Le véritable enjeu est inverse :

\begin{aiquote}
Que devient la société lorsque les capacités cognitives artificielles deviennent abondantes, accessibles et intégrées à presque toutes les activités humaines ?
\end{aiquote}

Ce livre part d'une hypothèse qu'il n'affirmera jamais comme une certitude :

\textbf{la révolution de l'IA ne réside peut-être pas dans le fait que les machines deviennent intelligentes, mais dans le fait que l'intelligence elle-même devient une ressource beaucoup plus largement accessible.}

Si cette hypothèse se vérifiait, elle aurait des répercussions en cascade sur les rapports qui, depuis des siècles, ont structuré nos sociétés : ceux entre intelligence, travail, valeur, pouvoir, organisation et société.

Il faut toutefois prendre la mesure d'une prudence élémentaire. Toute technologie est ambivalente. Elle ouvre des possibles, elle n'impose pas un destin. Aussi, plutôt que de décrire ce qui \emph{adviendra}, ce livre se propose d'explorer ce qui \emph{devient possible}. Car l'après n'est pas tracé d'avance. Il se construit, peu à peu, par les choix que nous ferons. Pour comprendre le poids de ce basculement éventuel, il nous faut d'abord regarder en arrière. Il faut comprendre que nous n'avons rien inventé d'absolument nouveau en ayant l'idée d'externaliser une partie de notre intelligence.

C'est là que commence notre réflexion.
